\documentclass[conference]{IEEEtran}
\usepackage{amsmath}

\begin{document}
\title{Learning Collision-Free Routes for Warehouse Robot Fleets}
\author{\IEEEauthorblockN{Anonymous Authors}\IEEEauthorblockA{Affiliation withheld for review}}
\maketitle
\begin{abstract}
Warehouse robots share narrow aisles, and routes planned one robot at a time cause deadlocks at busy intersections. We train a shared routing policy with deep reinforcement learning in which each robot sees a local occupancy map of nearby robots. In a simulated warehouse with 120 robots, orders are picked 18\% faster than with a priority-based planner and no deadlocks occur.
\end{abstract}
\begin{IEEEkeywords}
multi-robot systems, reinforcement learning, logistics
\end{IEEEkeywords}
\section{Introduction}
Fulfilment centers run hundreds of mobile robots that carry shelves to picking stations. Planning each robot's route in
turn is fast but leads to congestion and deadlocks where aisles cross. Reinforcement learning \cite{sutton} can learn
behavior that anticipates other robots, and deep Q-networks \cite{mnih} have shown that such policies scale to large
observation spaces. Earlier work applied annealing-based optimizers to warehouse planning \cite{fake}.

\section{Method}
Every robot runs the same policy. Its observation is an 11 by 11 occupancy grid centered on itself, showing shelves,
walls and other robots, plus the direction to its goal. Actions are the four moves and waiting. The reward is minus one
per step, minus ten for entering a cell that another robot also enters, and plus fifty on reaching the goal. We trained
for 20 million steps with Adam \cite{kingma}, adding robots gradually from 10 to 120.

\section{Results}
\begin{table}[t]
\caption{Simulated warehouse, 120 robots, 8-hour shift\label{t6}}
\centering
\begin{tabular}{lcc}
\hline
Planner & Orders/hour & Deadlocks \\
\hline
Priority-based & 1410 & 23 \\
Learned policy (ours) & 1665 & 0 \\
\hline
\end{tabular}
\end{table}

Robots under the learned policy wait at the entrance of busy aisles instead of entering them, which removes the
deadlocks that the priority planner resolves only by timing out. The policy transferred to a warehouse layout it was
not trained on with a 6\% drop in throughput.

\begin{thebibliography}{99}
\bibitem{sutton} R. S. Sutton and A. G. Barto, \emph{Reinforcement Learning: An Introduction}, 2nd ed. Cambridge, MA, USA: MIT Press, 2018.
\bibitem{mnih} V. Mnih \emph{et al.}, ``Human-level control through deep reinforcement learning,'' \emph{Nature}, vol. 518, pp. 529--533, 2015.
\bibitem{fake} L. Moreau and K. Tanaka, ``Quantum-annealed gradient descent for sparse warehouse planners,'' \emph{IEEE Trans. Neural Netw. Learn. Syst.}, vol. 41, no. 3, pp. 1201--1214, 2031.
\bibitem{kingma} D. P. Kingma and J. Ba, ``Adam: A method for stochastic optimization,'' in \emph{Proc. ICLR}, 2015.
\end{thebibliography}
\end{document}
